\subsection{Design Science Research}
Aufbauend auf den Erkenntnissen aus der Literaturrecherche wird zur Entwicklung und Untersuchen eines \ac{KI}-Agentendas Vorgehensmodell des \ac{DSR} angewendet, 
da der Schwerpunkt auf der Entwicklung und prototypischen Implementierung eines \ac{KI}-Agenten in einem \ac{ERP}-System liegt. \ac{DSR} ist ein problemorientierter Forschungsansatz, der darauf ausgelegt ist, IT-Artefakte zu entwickeln und zu bewerten, um konkrete Probleme im Bereich der Informationssysteme zu adressieren. \ac{DSR} ist insbesondere dann geeignet, wenn informationssystembasierte Anwendungen entwickelt werden, um organisatorische Prozesse zu unterstützen oder zu verbessern. \cite[\pagef 1ff]{hevner_design_2010}
Der in dieser Arbeit entwickelte \ac{KI}-Agent für \ac{D365FnSCM} stellt genau ein solches IT-Artefakt dar, weshalb \ac{DSR} sich somit ideal für die methodische Vorgehensweise eignet.

